\def\level{BTS}  % Classe à droite
\def\course{CIEL 2}
\def\chaptername{Lois continues} % Titre au centre
\def\docpart{Loi normale~--~TP2} % Partie du cours
\def\docname{C227}        % Identifiant à gauche

\pagestyle{cours_FP}
\makeheaderBTS{} % Affiche le bandeau

\begin{exercice_cours}[\quad]{}

Un revendeur dematériel photographique désire s’implanter dans une galerie marchande. 
Il estime qu’il pourra vendre 40 appareils par jour et les ventes sont deux à deux indépendantes.
Une étude lui amontré que, parmi les différentes marques disponibles, la marque A réalise $38.6\%$
du marché.

On note $X$ la variable aléatoire qui, pour une journée donnée, compte le nombre d’appareils de marque A vendus par le revendeur.

\begin{enumerate}
	\item Déterminer la loi de probabilité de $X$.
	\item Calculer $P(X=10)$ et $P(X=20)$.
	\item Déterminer la probabilité que le revendeur vende au moins 15 appareils de marque A au cours d’une journée.
	\item Calculer l’espérance de $X$ et donner une interprétation de ce résultat.
\end{enumerate}

On décide d'approcher la loi de $X$ par une loi normale $N(\mu, \sigma)$.

\begin{enumerate}
	\item Expliquer pourquoi il est possible d'effectuer cette approximation.
	\item Déterminer les paramètres $\mu$ et $\sigma$ de la loi normale $N(\mu, \sigma)$.
	\item Calculer $P(X=10)$ et $P(X=20)$ en utilisant la loi normale $N(\mu, \sigma)$.
	\item Déterminer la probabilité que le revendeur vende entre 19 et 21 objets.
\end{enumerate}
\end{exercice_cours}

	\begin{exercice_cours}[\quad]{}

Le diamètre intérieur moyen d'un échantillon de 200 corps de stylos produits par une machine est de $0,502$ cm et l'écart-type
moyen est de $0,005$ cm. Ne peuvent être acceptées pour des opérations de montage automatiques de montage qui suivent que
les pièces dont le diamètre est compris entre $0,496$ et $0,508$ cm, les autres étant considérés comme défectueuses.
Quel est alors le pourcentage de corps de stylos défectueux, sachant que les diamètres des pièces sont distribués normalement ?
\end{exercice_cours}


\begin{exercice_cours}[\quad]{}
Une étude a été menée sur des sportifs
amateurs de course à pied de catégorie 3 (pratiquant 2
à 4 fois par semaine). D’après cette étude, la fréquence cardiaque
maximale (FCM) d’un tel sportif suit une loi normale
de paramètre $\mu = 188$ et $\sigma=6,95$.
Quelle est la probabilité qu’un tel sportif choisi au hasard ait
une FCM comprise entre 179 et 180 ?
\end{exercice_cours}

\begin{exercice_cours}[\quad]{}
Le taux d’hématocrite (en \%) est le volume
occupé par les globules rouges dans les sang. Ce taux
(pour les hommes) est supposé suivre la loi normale de paramètres
$\mu = 45,5$ et $\sigma = 3$. Il est considéré normale s’il se situe
dans l’intervalle $\ff{39,5}{51,5}$.

\begin{enumerate}
	\item Quelle est la probabilité d’avoir un taux d’hématocrite
normale ?
	\item Lors d’un contrôle antidopage, le taux d’hématocrite
du vainqueur du tour de France 1996 aurait été supérieur
à 60. Quelle est la probabilité d’avoir un taux d’hématocrite supérieur à 60 ?
\end{enumerate}

\end{exercice_cours}


\begin{exercice_cours}[\quad]{}
Soit X une variable aléatoire suivant un loi
binomiale $B(100; 0,2)$ . En utilisant une approximation de
cette loi par la loi normale dont on précisera les paramètres,
calculer une valeur approchée de $P(X = 20)$, $P(X \leq 20)$,
$P(186 \leq X \leq 22)$ et de $P(X > 18)$


\end{exercice_cours}

\begin{exercice_cours}[\quad]{}
	Une compagnie aérienne estime à $0,1$ la probabilité qu’un
client ayant réservé sa place ne se présente pas à l’embarquement.

Sur le vol MA 2012, l’avion a une capacité de 300 places. Pour
optimiser son remplissage, la compagnie a accepté plus de
300 réservations. Ce faisant, elle court le risque que se présentent
à l’embarquement plus de 300 personnes ayant réservé,
auquel cas elle devra indemniser ceux qui ne pourront
embarquer.

On note $n$ le nombre de réservations acceptées par la compagnie,
et $X$ la variables aléatoire indiquant le nombre de
personnes ayant réservé qui se présentent à l’embarquement.
On suppose les comportements des clients indépendants
les uns des autres

\begin{enumerate}
	\item  Justifier que $X$ suit une loi binomiale dont vous préciserez
les paramètres.
	\item Justifier que cette loi binomiale peut être approché
par une loi normale dont vous préciserez les les paramètres.
	\item On suppose dans cette question que $n = 324$.
Quelle probabilité que la compagnie ne puisse pas embarquer
tous les passagers qui se présentent ?
	\item La compagnie souhaite limiter à $1$ \% le risque de
ne pouvoir embarquer tous les passagers qui se présentent.
Déterminer le nombre maximum de places qu’elle
peut proposer à la réservation.
\end{enumerate}
\end{exercice_cours}

\begin{exercice_cours}[\quad]{}

	Soit $X$ une variable aléatoire suivant la loi normale de paramètres $\mu=50$ et $\sigma=5$. Déterminer le réel $a$ arrondi à $10^{-2}$ près vérifiant:
\begin{enumerate}
	\item $P(X\leqslant a)=0,2$,
	\item $P(X\geqslant a)=0,9$,
	\item $P(50-a \leqslant X \leqslant 50+a)=0,95$.
\end{enumerate}
\end{exercice_cours}

\begin{exercice_cours}[\quad]{}

	Une machine fabrique des résistors. La variable aléatoire $X$ associe à chaque résistor sa résistance exprimée en ohms. $X$ suit la loi normale d'espérance 80 et d'écart type 2.
On prélève un résistor au hasard.

\begin{enumerate}
	\item Calculer la probabilité que sa résistance soit inférieure à $81,5$.
	\item Un résistor est conforme si sa résistance est comprise entre $79,85$ et $81,15$. Quelle est la probabilité que le résistor soit conforme ?
	\item Déterminer le réel $h$ tel que $95 \%$ des résistors aient une résistance comprise entre $80-h$ et $80+h$.
\end{enumerate}
\end{exercice_cours}

\begin{exercice_cours}[\quad]{}

On s'intéresse à un chantier de construction d'un tronçon de TGV. Les travaux de terrassement nécessitent la mise à disposition d'une flotte importante de pelles sur chenilles et de camions-benne. La réalisation de l'ouvrage nécessite de grandes quantités de béton.
Les résultats approchés seront arrondis à $10^{-3}$.

\medskip

\textbf{A. Loi normale}

\medskip

On note $X$ la variable aléatoire qui, à chaque pelle prélevée au hasard dans la flotte, associe le nombre de m3 de matériaux extraits pendant la première heure de chantier.
On suppose que $X$ suit la loi normale d'espérance 120 et d'écart type 10.

\begin{enumerate}
	\item Calculer $P(110 \leq X \leq 130)$.
	\item  Calculer la probabilité que la pelle prélevée extrait moins de 100 m3 pendant sa première heure de chantier.
\end{enumerate}

\medskip

\textbf{B. loi binomiale}
\medskip

On note $\mathcal{E}$ l'événement "un camion-benne pris au hasard dans la flotte n'a pas de panne ou de sinistre pendant le premier mois de chantier". On suppose que la probabilité de $\mathcal{E}$ est $0,9$.

On prélève au hasard 10 camions-benne dans la flotte pour les affecter à une zone de chantier. Le nombre de camion benne de la flotte est assez important pour que l'on puisse assimiler ce prélèvement à un tirage avec remise de 10 camions-benne.

On note $Y$ la variable aléatoire, à tout prélèvement de ce type, associe le nombre de camions-benne n'ayant pas eu de panne ou de sinistre pendant le premier mois de chantier.

\begin{enumerate}
	\item Justifier que $Y$ suit une loi binomiale et on précisera les paramètres.
	\item Calculer la probabilité que, dans un tel prélèvement, aucun des 10 camions-benne n'aient de panne ou de sinistre pendant le premier mois de chantier.
\end{enumerate}
	\end{exercice_cours}

	\begin{exercice_cours}[\quad]{}

\emph{On arrondira les probabilités à $10^{-3}$.}

Dans ce problème, on s'intéresse à une production de pots de confiture dans une usine.

On considère un événement $\mathcal{E}$ : "un pot à une masse inférieure à 490 grammes". Une étude a permis d'admettre que la probabilité de cet événement est de $0,2$.


\medskip

\textbf{A. Loi normale}

\medskip

On prélève au hasard 10 pots dans la production totale.

On suppose que le nombre de pots est assez important pour que l'on puisse assimiler ce prélèvement à un tirage avec remise de 10 pots. On considère la variable aléatoire $X$ qui, à tout prélèvement de 10 pots, associe le nombre de pots dont la masse est inférieure à 490 grammes.

	\begin{enumerate}
		\item Expliquer pourquoi $X$ suit une loi binomiale. En préciser les paramètres.
		\item Calculer la probabilité de l'événement $A$ : "parmi les 10 pots, il y a exactement 2 pots dont la masse est inférieure à 490 grammes".
		\item Calculer la probabilité de l'événement $B$ : "parmi les 10 pots, il y a plus de 2 pots dont la masse est inférieure à 490 grammes".
	\end{enumerate}


\medskip

\textbf{B. loi binomiale}
\medskip

On prélève au hasard 100 pots dans la production totale. On considère la variable aléatoire $Y$ qui, à tout prélèvement de 100 pots, associe le nombre de pots dont la masse est inférieure à 490 grammes.

\begin{enumerate}
	\item On décide d'approcher la loi de la variable aléatoire $Y$ par une loi normale. Utiliser cette approximation pour calculer la probabilité que plus de 25 pots ont une masse inférieure à 490 grammes.
	\item Calculer de même la probabilité que moins de 18 pots ont une masse inférieure à 490 grammes.

\end{enumerate}

	\end{exercice_cours}
